\chapter{The paired-prime groups}
\label{ch:construction}
\chapterintro{A binary selector chooses one prime from each pair. Its indicator determines a rational line in the ambient quotient, and an explicit finite cocycle turns that line into a countable group.}

\section{Choosing a rational line}
Let $A\subseteq\N$, write $a_n=\one_A(n)$, and define
\[
S_A=\{p_{2n+a_n}:n\in\N\}.
\]
Exactly one prime from each pair is selected. Both $S_A$ and its complement are infinite. Let $u_A\in P$ be its indicator:
\[
(u_A)_p=\begin{cases}1,&p\in S_A,\\0,&p\notin S_A.\end{cases}
\]
Because $u_A$ has infinite support, $q(u_A)\ne0$. We can therefore define a rational line and its preimage by
\[
W_A=\Q q(u_A),\qquad G_A=q^{-1}(W_A).
\]
This definition produces an actual group for every set $A$. It does not claim an effective method for obtaining the selector bits of an arbitrary $A$.

\begin{figure}[tbp]
\centering\includegraphics[width=\textwidth]{\figfile{02_prime_pairs.pdf}}
\caption{Each bit selects exactly one prime from its pair. A filled point is selected. The torsion factor $C_p$ is present in $G_A$ whether or not $p$ is selected. Selection controls the infinite-order lift and its divisibility, not the presence of the torsion factor.}
\label{fig:primepairs}
\end{figure}

\section{A section with finite arithmetic defects}
For a reduced rational $r=a/b$, with $b>0$, define
\begin{equation}
\label{eq:section}
(s_A(a/b))_p=
\begin{cases}
a b^{-1}\pmod p,&p\in S_A\text{ and }p\nmid b,\\
0,&p\notin S_A\text{ or }p\mid b.
\end{cases}
\end{equation}
The inverse is used only when the denominator is a unit in $\F_p$. The convention at denominator primes is essential. It is a choice of representatives, not a claim that division by $p$ is defined in $C_p$.

The difference $b s_A(a/b)-a u_A$ is supported on the prime divisors of $b$. It follows that
\[
q(s_A(r))=r q(u_A).
\]
Thus every element of $G_A$ is $s_A(r)+d$ for some $r\in\Q$ and $d\in D$.

\begin{theorem}[Unique normal form]
\label{thm:normalform}
Every element of $G_A$ has a unique expression
\[
g=s_A(r)+d,\qquad (r,d)\in\Q\times D.
\]
Consequently $G_A$ is countably infinite, with a countable normal-form domain that does not depend on the selector.
\end{theorem}
\begin{proof}
Existence follows from the preimage definition and the displayed quotient identity. If two forms are equal, applying $q$ gives $(r-r')q(u_A)=0$. The vector $q(u_A)$ is nonzero in a rational vector space, so $r=r'$. The finite-support terms are then equal. Both $\Q$ and $D$ are countable, while the quotient is the infinite group $\Q$.
\end{proof}
The domain can be named without knowing $A$. The group operation on those names will contain the selector information.

\section{The cocycle law}
Define
\[
c_A(r,t)=s_A(r)+s_A(t)-s_A(r+t).
\]
The sum is taken in $P$, but its value belongs to $D$.

\begin{lemma}[A finite support bound]
\label{thm:cocyclesupport}
For all $r,t\in\Q$,
\[
\supp c_A(r,t)\subseteq
\{p:p\mid\den(r)\den(t)\}.
\]
\end{lemma}
\begin{proof}
At a prime not dividing either denominator, all three rationals have well-defined reductions. On a selected coordinate, reduction respects addition, so the defect is zero. On an unselected coordinate, all section values are zero. The reduced denominator of $r+t$ has no prime divisor absent from the two original denominators.
\end{proof}

\begin{theorem}[Explicit group law]
\label{thm:cocyclelaw}
In normal-form coordinates, addition is
\begin{equation}
\label{eq:grouplaw}
(r,d)+(t,e)=(r+t,d+e+c_A(r,t)).
\end{equation}
The identity is $(0,0)$ and the inverse of $(r,d)$ is $(-r,-d)$. The law is associative and commutative, and the map $(r,d)\mapsto s_A(r)+d$ is an isomorphism onto the subgroup $G_A\le P$.
\end{theorem}
\begin{proof}
Substitution into addition in $P$ gives the formula. The section satisfies $s_A(0)=0$ and $s_A(-r)=-s_A(r)$, so $c_A(r,0)=0$ and $c_A(r,-r)=0$. These identities give the stated identity and inverse.

Symmetry is immediate. The cocycle equation
\[
c_A(r,t)+c_A(r+t,v)=c_A(t,v)+c_A(r,t+v)
\]
follows by expanding and cancelling the section terms. Expanding the two associations in \eqref{eq:grouplaw} now proves associativity. Uniqueness of normal forms makes the ambient map bijective, and the addition formula makes it a homomorphism.
\end{proof}
This is the standard symmetric cocycle construction of an abelian extension \cite[Section 3.2]{Lupini}. Here the particular section makes each individual correction finitely supported and exactly computable from finitely many selector queries.

\begin{proposition}[One cocycle value reads one coordinate]
\label{thm:cocyclebit}
For every prime $p$,
\[
c_A(1/p,(p-1)/p)=-\one_{S_A}(p)e_p.
\]
\end{proposition}
\begin{proof}
The support bound leaves only the coordinate $p$. At that coordinate both fractional lifts are zero, and the lift of their sum, which is one, is $\one_{S_A}(p)$.
\end{proof}
Dropping the cocycle is therefore not a harmless simplification. The direct product operation on $\Q\times D$ removes exactly these defects and gives a different, split extension.

\section{Structure of the extension}
By Chapter~\ref{ch:purity}, the group is pure in $P$, reduced, residually finite, and dense. Its torsion is $D$, and the rational coordinate gives the exact sequence
\begin{equation}
\label{eq:extension}
0\longrightarrow D\longrightarrow G_A\xrightarrow{\rho_A}\Q\longrightarrow0,
\qquad \rho_A(s_A(r)+d)=r.
\end{equation}
It does not split.

There is an equivalent cocycle obstruction. Suppose a function $b:\Q\to D$ modified the section to an additive map $r\mapsto s_A(r)+b(r)$. Then
\[
c_A(r,t)+b(r)+b(t)-b(r+t)=0
\]
for every $r,t$. The modified section would split \eqref{eq:extension}. Hence no such function exists. The cocycle is the defect of a section into $P$, but it is not removable by a correction valued in $D$.

\begin{proposition}[Finitely generated subgroups]
\label{thm:fgsubgroups}
Every finitely generated subgroup of $G_A$ is either finite cyclic or isomorphic to $\Z\oplus C_m$ for a squarefree positive integer $m$, with $C_1$ trivial.
\end{proposition}
\begin{proof}
Let $L\le G_A$ be generated by finitely many elements. Its image in $\Q$ is a finitely generated subgroup of $\Q$, hence either zero or infinite cyclic. If it is zero, then $L\le D$ and its finite support places it inside a finite cyclic group.

Otherwise choose a lift $x\in L$ of a generator of its rational image. Subtracting suitable integer multiples of $x$ from the given generators produces finitely many elements of $D$. They generate the kernel of $L\to\rho_A(L)$, which is finite cyclic of squarefree order. The subgroup $\langle x\rangle$ intersects that torsion kernel trivially, giving the direct sum.
\end{proof}
The extensions over finitely generated rational images split, but the extension over all of $\Q$ does not. This is another exact distinction between finite algebraic stages and the global realization.
